\begin{abstract}
We prove the metric Blaschke conjecture: every connected closed smooth
Riemannian manifold of positive dimension whose global injectivity radius
equals its diameter is, up to scale, a standard compact rank-one symmetric
space.
\end{abstract}

\section{Introduction}
\label{sec:introduction}

On a round sphere, every geodesic minimizes from its starting point to
the antipode, at a distance equal to the diameter. The standard real,
complex and quaternionic projective spaces and the Cayley plane have
the same common minimizing-distance property. The metric Blaschke
conjecture asks whether this property characterizes these compact
rank-one symmetric spaces.

For a connected closed smooth Riemannian manifold $(M,g)$, write
\[
 \inj(M,g)=\inf_{p\in M}\inj_p(M,g),
 \qquad
 \diam(M,g)=\sup_{p,q\in M}d_g(p,q).
\]
Here $\inj_p(M,g)$ is the supremum of the radii on whose open tangent
balls $\exp_p$ is a diffeomorphism onto its image. We call $M$ a
\emph{Blaschke manifold} if it has positive dimension and
$\inj(M,g)=\diam(M,g)$. The global equality implies that every
unit-speed geodesic minimizes precisely up to the common diameter;
we give the short deduction in Lemma~\ref{geom:cut-time}.

\begin{theorem}[The metric Blaschke theorem]\label{thm:blaschke}
Let $(M,g)$ be a connected closed smooth Riemannian manifold of
positive dimension. If $\inj(M,g)=\diam(M,g)$, then $(M,g)$ is
isometric, up to multiplication of the metric by a positive constant,
to a round sphere, a standard real, complex or quaternionic projective
space, or the Cayley plane.
\end{theorem}

The classification includes every quaternionic dimension. Its hypothesis
is the global equality just defined. Classical Blaschke structure
associates to $M$ a compact rank-one symmetric \emph{model} $M_0$ by
its cut-locus and integral cohomology type. Throughout the volume
comparisons, the standard metric $g_0$ on $M_0$ is scaled to have
the same diameter as $(M,g)$. Prime geodesics have length twice
the diameter, so fixing the diameter is equivalent to fixing
their prime period.

\subsection{The problem and its predecessors}

Blaschke's original question concerned surfaces with a common
reconvergence property for geodesics. Green settled the surface
problem \citep{Blaschke1921,Green1963}; the later geometric formulation
and its development are described in \citet{Besse1978} and in the
corrected account of \citet{McKay2016}. The classical theory connects
the minimizing condition to closed geodesics, smooth cut loci and
fibrations of tangent spheres by great spheres. The topological
restrictions of Bott and Samelson explain the appearance of the
compact rank-one symmetric models
\citep[Theorem~7.2]{Besse1978}.

The sphere and real-projective metric cases were established through
the work of Berger, Kazdan, Weinstein and Yang; see
\citet[Appendices~D--E]{Besse1978} and
\citet[Sections~11--12]{McKay2016}. Berger's sphere argument combines
the geometric volume problem with Kazdan's scalar integral inequality
\citep{Berger1978AppendixD,Kazdan1978AppendixE}. Berger subsequently
proved the sharp isoembolic inequality by a weighted comparison of
Jacobi determinants and Jensen's inequality
\citep[pp.~263--264]{Berger1980}. These comparisons are central
predecessors of the analytic part of the present proof.

Under additional hypotheses on families of totally geodesic round
spheres of curvature $4$, \citet[Theorem~26, pp.~8--11]
{RovenskiiToponogov1996} obtain a sharp projective volume comparison.
Their Jacobi determinant splits into an exact curvature-$4$ factor
and a complementary factor controlled by Kazdan's inequality.
Our fixed-chart Schur-complement argument uses the cut-time
intersection pattern without requiring that geometric splitting.

For the projective models, it is useful to distinguish three kinds
of information: topology, volume at a fixed geodesic period, and
metric isometry. Yang proved that a Blaschke manifold of complex
projective \emph{homotopy type} has the model volume at the same
prime geodesic period \citep[Theorem, p.~380]{Yang1991}. For the
projective planes, smooth incidence geometry and the classification
of \citet[Corollary~C]{KramerStolz2007} identify the point manifold
diffeomorphically. Wilking records the equal-volume theorem for
a manifold homeomorphic to a compact rank-one symmetric space whose
unit geodesics all have the same least period, attributing it to
Reznikov \citep[p.~282]{Wilking2001}. Reznikov's earlier work
on the quaternionic weak Blaschke conjecture
\citep{Reznikov1985HP} is another predecessor of this volume
approach. The homeomorphism hypothesis in the result recorded
by Wilking, and the other hypotheses just stated, will be
verified at their uses. In higher quaternionic dimensions our upper
bound is intrinsic: it is proved from the Blaschke geometry and
integral topology in Sections~\ref{inc:section}--\ref{alg:section}.

\subsection{Two volume bounds and their common equality case}

The first part of the proof establishes a comparison valid for every
model type.

\begin{theorem}[Sharp volume comparison]\label{thm:comparison}
Every Blaschke manifold and its model at the same diameter satisfy
\[
 \Vol(M,g)\ge\Vol(M_0,g_0).
\]
Equality holds if and only if $(M,g)$ is isometric to $(M_0,g_0)$.
\end{theorem}

\begin{corollary}[Reduction to an upper volume bound]
\label{thm:volume-reduction}
If a Blaschke manifold has volume at most that of its model at the
same diameter, then it is isometric to that model.
\end{corollary}

After the circle, sphere and real-projective cases are separated,
the analytic comparison in each remaining type starts along one
closed geodesic. Normal Jacobi fields that vanish at a specified
time form a Lagrangian plane in the symplectic space of Jacobi
fields. A fixed graph chart separates the determinant of their
endpoint pairing into a regular factor and a Schur-complement
factor. The size of the Schur block is the dimension of the space
of normal Jacobi fields vanishing at both the starting point and
its cut point; the regular block has the complementary normal
dimension. Both determinant factors can be bounded by scalar
interval integrals, including when the actual endpoint order is
reversed. This construction uses the cut-time intersection pattern;
it does not require parallel eigenspaces of the curvature operator.

The resulting scalar estimates are combined by a logarithmic form
of Kazdan's weighted cancellation
\citep[Appendix~E, especially E.4--E.8]{Kazdan1978AppendixE}.
Its reflection symmetry cancels the integrated endpoint terms
against interval averages. The scalar reference profiles are
reciprocal-square sines, one on the full geodesic period and the
other on a half-period. Folding the latter factor onto the two
half-periods produces four weights, one for each choice of a
half-period for each endpoint. Their equal, symmetric row and
column sums allow both comparisons to contribute with the
correct multiplicities. A nonnegative strict-convexity remainder
records equality. Integrating over the unit tangent bundle gives
the volume lower bound. Equality makes every radial volume density
equal to its model value, and Szab\'o's classification of compact
simply connected harmonic manifolds then gives metric rigidity
\citep{Szabo1990,Szabo1988}.

The quaternionic upper bound has a different source. Fix dimension
$4r$, with $r\ge2$. Pairs of points at the cut distance form a
smooth incidence manifold. Its actual normal four-plane bundle
identifies the great-three-sphere families of tangent directions.
Two different direction planes at one point intersect only at zero.
The rank of their complement forces a formula for the tangent
Pontryagin class; the signature theorem and a uniform estimate in
every rank then determine its coefficient. Characteristic-class
and signature theory enter here in their usual integral and
oriented forms \citep{MilnorStasheff1974}.

Next consider the space of oriented geodesics, forgetting the
marked starting point. The period-normalized contact form expresses
volume as a top power of the Euler class of its geodesic-circle
bundle. Integral Gysin sequences first determine the lattice of
classes and the determinant of a middle pairing. The first negative
bundle of the free-loop energy transfers the incidence coefficient
to this space. This step uses the
fixed-endpoint Morse index theorem of \citet{Milnor1963} and the
equivariant negative-bundle attachment developed in closed-geodesic
Morse theory; we use the formulation of
\citet[Section~7]{OttosenBokstedt2007}, following
\citet{Klingenberg1978}. The required primitive attaching coefficient
is established over every prime field. The attaching map from
the sphere bundle of the descended negative bundle to $M$ also
pulls back its cohomology generator; the vanishing of the
$(r+1)$st power supplies the relations for the final algebraic
calculation. Finally, a complete-intersection calculation turns
the middle determinant
into the sharp upper volume bound. Rational cohomology
alone would not retain the arithmetic that normalizes this bound.

Section~\ref{geom:section} fixes the geometric conventions.
Sections~\ref{sec:lagrangian} and~\ref{weighted:section} prove the
geodesic comparison, and Section~\ref{rig:section} completes its
global volume and equality argument. Sections~\ref{inc:section}
and~\ref{sig:section} establish the quaternionic characteristic
class. Sections~\ref{geo:section} and~\ref{alg:section} convert it
to the upper bound. Section~\ref{completion:section} checks the
remaining classical inputs and completes the metric classification.

\subsection{Recent related rigidity results}

There are also rigidity theorems with additional geometric hypotheses.
Shankar, Spatzier and Wilking treat an upper sectional-curvature bound
\citep[Proposition~2.1]{ShankarSpatzierWilking2005}, and Su, Sun and
Wang prove the corresponding normalized lower-curvature result
\citep{SuSunWang2016}. Connell, Islam, Nguyen and Spatzier obtain
rigidity under curvature bounds and agreement with a model away
from a prescribed convex set \citep{ConnellIslamNguyenSpatzier2026}.

Song's September 2026 preprint proves Blaschke metric rigidity
assuming an adapted complex structure on the entire tangent bundle
\citep[Theorem~1.3]{Song2026}. This entire-extension hypothesis is
additional to the Blaschke condition.
